\chapter{Leitura no contexto}

\begin{objetivo}
Estudar os textos centrais um a um: o contexto, o que o texto diz e como
contribui para responder à pergunta.
\end{objetivo}

\begin{fichatexto}[Sangue e vida]{Lv 17:11}
\begin{campo}{Contexto}
\rb{Lv 17} regula o abate de animais: todo sacrifício deve ser levado à
tenda do encontro (vv.~1–9), e o sangue não deve ser comido, nem pelo
israelita nem pelo estrangeiro que mora entre eles (vv.~10–16). O v.~11 dá o
motivo da proibição do v.~10.
\end{campo}

\interlinear[titulo={Lv 17:11 · WLC}, legenda=false]{dados/biblia/lv017_11.hbo.tsv}

\begin{campo}{Observações sobre o texto}
\hb{נֶפֶשׁ} (\tl{nep̄eš}) aparece três vezes: a \enquote{vida} da carne está no
sangue; o sangue faz expiação pelas \enquote{vidas} (\tl{nap̄šōtêkem}) de
Israel; e \tl{bannep̄eš}. Neste último caso, a preposição \hb{בְּ} pode ser
lida como instrumental (\enquote{por meio da vida}) ou de troca
(\enquote{em troca da vida}); as traduções se dividem. O verbo \hb{כָּפַר}
(Piel: \enquote{fazer expiação}) aparece duas vezes.
\end{campo}

\campovazio[3]{Contribuição para o tema}{}
\end{fichatexto}

\begin{fichatexto}[Não comer sangue]{At 15:28, 29}
\begin{campo}{Contexto}
A carta da reunião de apóstolos e anciãos em Jerusalém às congregações de
cristãos não judeus (\rb{At 15:22–29}), depois da discussão sobre a
circuncisão (\rb{15:1–21}). Os mesmos quatro itens aparecem em \rb{15:20} e
são lembrados em \rb{21:25}.
\end{campo}

\interlinear[titulo={At 15:28–29 · SBLGNT}, legenda=false]{dados/biblia/at015_28-29.grc.tsv}

\begin{campo}{Observações sobre o texto}
\gr{ἀπέχεσθαι} é o mesmo verbo de \rb{Lc 6:24} (\gr{ἀπέχετε}), mas na voz
média, com genitivo: \enquote{abster-se de}, \enquote{manter-se longe de}. Na
ativa, em Lucas, era \enquote{receber por inteiro}. \gr{πλήν} aqui é
preposição (\enquote{exceto}), não a conjunção adversativa de \rb{Lc 6:24}.
Há variantes: o texto bizantino lê \gr{πνικτοῦ} (singular), e alguns
manuscritos do chamado texto ocidental omitem \enquote{estrangulado} e
acrescentam uma forma negativa da regra de ouro \parencite{metzger1994}.
\end{campo}

\campovazio[3]{Contribuição para o tema}{}
\end{fichatexto}

\begin{fichatexto}[Sangue e vida]{Gn 9:4}
\campovazio[3]{Contexto}{a quem Deus fala, em que momento da história}
\campovazio[3]{O que o texto diz}{com as minhas palavras}
\campovazio[3]{Contribuição para o tema}{}
\end{fichatexto}
